% ダイオード 課題（ワークシート，学生配布用）
% 解答例版は 03_diode_worksheet_answer.tex をコンパイルしてください。
\documentclass[student]{experimentguide}

\usepackage{tikz}
\usepackage{multirow}

\coursename{電気電子工学実験 I}
\graphicspath{{contents/03_diode/figures/}}

\makeatletter
% 解答（instructor オプション時のみ表示）
\ifeg@instructor
  \newcommand{\answer}[1]{\textcolor{EGWarn}{#1}}
  \newcommand{\ifanswer}[2]{#1}
\else
  \newcommand{\answer}[1]{}
  \newcommand{\ifanswer}[2]{#2}
\fi
\makeatother

% オシロスコープ画面の記入欄（横10 DIV × 縦8 DIV，1 DIV = 5 mm）
\newcommand{\scopegrid}{%
  \begin{tikzpicture}
    \draw[very thin,black!45,step=5mm] (0,0) grid (50mm,40mm);
    \draw[black!75] (0,20mm) -- (50mm,20mm);
    \draw[black!75] (25mm,0) -- (25mm,40mm);
    \draw[thick] (0,0) rectangle (50mm,40mm);
  \end{tikzpicture}%
}

\newcommand{\wavebox}[2]{% #1: 解答写真, #2: 見出し
  \begin{minipage}[t]{0.32\linewidth}
    \centering
    \ifanswer{\includegraphics[width=50mm,height=40mm,keepaspectratio]{#1}}{\scopegrid}\par
    \vspace{0.3\baselineskip}
    \small #2
  \end{minipage}%
}

\pagestyle{empty}

\begin{document}
% ============================================================
% ダイオード 課題（ワークシート）本文
% \answer{...} の中身は解答例版（03_diode_worksheet_answer.tex）だけに表示されます。
% ============================================================
\begin{center}
  {\sffamily\gtfamily\bfseries\Large 電気電子工学実験I　ダイオード　記録用紙%
    \ifanswer{（解答例）}{}}
\end{center}

\vspace{0.6\baselineskip}
\begin{flushright}
  学籍番号：\underline{\hspace{12\zw}}\quad
  氏名：\underline{\hspace{12\zw}}
\end{flushright}
\vspace{0.8\baselineskip}

\section*{課題1}

実験1及び実験2より表1を完成させてください。

\begin{table}[h]
  \centering
  \caption*{表1：各色のLEDの電圧と電流の関係}
  \renewcommand{\arraystretch}{1.5}
  \begin{tabular}{c>{\centering\arraybackslash}p{9\zw}>{\centering\arraybackslash}p{9\zw}}
    \toprule
    LEDの色 & 電圧 $V$（V） & 電流 $I$（A） \\
    \midrule
    赤色   & \answer{1.80} & \multirow{4}{*}{0.001} \\
    黄緑色 & \answer{1.90} & \\
    青色   & \answer{2.55} & \\
    白色   & \answer{2.60} & \\
    \bottomrule
  \end{tabular}
\end{table}

\section*{課題2}

実験3より(a)交流電源から出力される交流電圧波形，(b)整流回路にLEDを接続したときの交流電圧波形，(c)整流回路にコンデンサを接続したときの交流電圧波形をそれぞれ図示してください。

\begin{center}
\wavebox{diode_answer_wave_a.jpg}{図(a)\quad 交流電源から出力される交流電圧波形}\hfill
\wavebox{diode_answer_wave_b.jpg}{図(b)\quad 整流回路にLEDと抵抗を並列接続したときの交流電圧波形}\hfill
\wavebox{diode_answer_wave_c.jpg}{図(c)\quad 整流回路にコンデンサと抵抗を並列接続したときの交流電圧波形}
\end{center}

\end{document}
